\pdfoutput=1

\section{Related work}
\label{related_work}

% Researchers have explored several classes of methods to enhance LLM reasoning. Single-agent approach falls into two categories: search-based \citep{brown2020language,graves2012sequence,snell2024scaling,wang2022self,lightman2023let} and post-training methods \citep{qin2024o1replicationjourneystrategic, ouyang2022training, hui2024qwen2, liu2024deepseek,wang2024openr, zhang2024llamaberrypairwiseoptimizationo1like, r1, xu2025towards}. Multi-agent approach includes aggregating static knowledge from off-the-shelf LLMs \citep{zhang2025if,chen2024s,du2023improving,liang2023encouraging,hong2023metagpt,zhang2024aflow} and improving reasoning ability via training \citep{ma2023red,kirchner2024prover}. Hierarchical reasoning approach further partitions reasoning into hierarchical processes \citep{saha2025learningplanreason,DBLP:journals/corr/abs-2410-03864,zhao2024marcoo1openreasoningmodels}. 
% RL training  has been applied directly to reasoning. Many works examine and improve RL training on LLMs \citep{DBLP:journals/corr/abs-2503-20783,yue2025vapo,yu2025dapo,zeng2025simplerlzoo,zhao2025echo,DBLP:journals/corr/abs-2405-15821}. Some works push RL training for single-turn to multi-turn \citep{jin2025searchr1,zhou2024archer,wang2025ragen}.
% See Appendix~\ref{sec:appx.related_work} for a detailed analysis of each line of work.

Drawing from the bitter lesson \citep{sutton2019bitter}, two methods that appear to scale effectively are searching and learning, aligning with current trends in large language models \citep{xu2025towards}. At present, researchers are leveraging these methods to maximize the capabilities of individual transformers, while other efforts are exploring architectures that involve multiple interacting entities. In this paper, we examine this divergence within the context of LLM reasoning, a capability that allows large language models to solve problems through logical reasoning, step-by-step analysis, and inference \citep{wang2024openr}.

\subsection{Single LLM Reasoning}
\label{related_work:single_llm}
Main research works in reasoning involving a single LLM utilize search-based and post-training methods. The fundamental elements of searching methods are text generation and evaluation. Generation schemes include In-Context Learning \citep{brown2020language}, Beam Search \citep{graves2012sequence}, and various tree-based searching \citep{snell2024scaling}; Evaluation approaches often use outcome accuracy, self-consistency \citep{wang2022self}, or process reward signal \citep{lightman2023let} as the criteria to select high-quality responses from the generated texts. Post-training method is another research line in opposition to pre-training. Popular training pipelines often involve specific data construction followed by Supervised Fine-tuning \citep{qin2024o1replicationjourneystrategic, ouyang2022training, hui2024qwen2, liu2024deepseek}, or reinforcement learning to interactively explore learning patterns \citep{wang2024openr, zhang2024llamaberrypairwiseoptimizationo1like, r1, xu2025towards}.

\subsection{Multiple LLM Reasoning}
\label{related_work:multi_llm}
Integrating multiple entities can potentially surpass the intelligence of the individual model \citep{chen2023agentverse}. With the rapid emergence of large language models showing a varying level of abilities, some studies have explored facilitating discussions among multiple off-the-shelf LLMs \citep{zhang2025if,chen2024s,wang2023unleashing,du2023improving,zhuge2023mindstorms,tang2023medagents,hao2025chatllm,akata2023playing,hong2023metagpt,zhang2024aflow}, taking the form of free discussion \citep{du2023improving,liang2023encouraging} or structured role assignments \citep{hong2023metagpt,zhang2024aflow}. Some have applied routing mechanisms to assign tasks to the most suitable expert models \citep{hu2024routerbench,stripelis2024tensoropera,ding2024hybrid,yue2025masrouter,chen2024routerdc} or merging mechanisms to develop more versatile models \citep{yadav2024matters,yu2024language,zhang2025nature}. Beyond aggregating static knowledge from multiple agents, multi-agent LLM training can also enhance reasoning capabilities. For example, multi-agent debates can generate diverse synthetic data, which can subsequently be used for supervised fine-tuning \citep{estornell2024accdebateactorcriticapproachmultiagent,li2024can,motwani2024maltimprovingreasoningmultiagent,dong2025stpselfplayllmtheorem,perez2022red,ye2025emergence,subramaniam2025multiagentfinetuningselfimprovement}. Reinforcement learning (RL) methods have also been adopted to improve LLM reasoning in areas such as alignment \citep{perez2022red,ma2023red} and legibility \citep{kirchner2024prover}. \cite{motwani2024maltimprovingreasoningmultiagent} utilize a three-agent system for generation and fine-tune the models using Direct Preference Optimization (DPO). Reinforcement Learning with Generative Reward Models (GenRM) \citep{mahan2024generative,ye2025iterative,jiao2024preference,wang2024self} represents another common approach of multi-agent training, where the reward signal is derived from the token probabilities of another LLM, coupled with the reasoning process. While our work aligns with these efforts, it diverges by using an additional tunable LLM to provide metacognitive instructions, guiding the low-level LLM during learning, rather than relying on a static GenRM. The most closely related works to ours are MAPoRL \citep{park2025maporlmultiagentpostcotrainingcollaborative} and COPRY \citep{DBLP:journals/corr/abs-2410-06101}. MAPoRL is a multi-agent debating framework that uses multi-agent reinforcement learning (MARL) with a learned verifier to fine-tune each LLM agent. COPRY duplicates an LLM into two agents, training them simultaneously in the roles of pioneer and observer using RL. 
\citet{shen2025satorireinforcementlearningchainofactionthought} trained with a novel Chain-of-Action-Thought (COAT) framework that embeds meta-action tokens for self-reflection and exploration into an autoregressive search process.
However, unlike our approach, which explicitly separates metacognition from plan execution, these methods do not decompose the reasoning process but instead focus on improving direct chain-of-thought generation. Furthermore, our experiments are conducted on a larger scale and include more challenging problems.

\subsection{Hierarchical Reasoning}\label{sec:ref3}
\label{related_work:hrl_reason}
Partitioning reasoning into hierarchical processes has been explored in prior research to make biological sense \citep{ye2018survey,langley2004hierarchical}. In the context of language models, a hierarchical structure has been used to facilitate diverse reasoning patterns, including planning \citep{puerta2025roadmap,sun2024retrieval,song2023llm,rana2023sayplan,chen2024autotamp,yan2023ask,xiao2024cellagent}, validation \citep{haji2024improvingllmreasoningmultiagent,xi2024enhancing} and self-refinement \citep{madaan2023self,kumar2024training,welleck2022generating}. For instance, EvalPlanner \citep{saha2025learningplanreason} is a framework that conducts reasoning through plan generation and execution. DOTS \citep{DBLP:journals/corr/abs-2410-03864} extends decomposition by integrating a tree-based searching method with Analysis, Solution, and Verification layers. Marco-o1 \citep{zhao2024marcoo1openreasoningmodels} focuses on open-ended problem-solving and abstract thinking, dynamically adjusting reasoning granularity and incorporating reflection mechanisms to enhance reasoning performance. Beyond these approaches, metacognition \citep{flavell1979metacognition} has been identified as another critical component of reasoning, referring to the intuitive understanding of one's own cognitive and reasoning processes \citep{gao2024meta,wang2023metacognitive}. \cite{wang2023metacognitive} proposed a metacognitive prompting strategy to improve large language model (LLM) capabilities. \cite{didolkar2024metacognitive} further developed a prompt-guided method that enables models to label math problems with the required skills and subsequently use these labels to solve new problems. \cite{gao2024meta} introduce meta-reasoner which use contextual multi-arm bandit to learn a high-level ``advisor'' over low-level reasoning process. \cite{xiang2025towards} provides a Meta-CoT framework to think about its own thinking. They use construction-based methods as well as reinforcement learning to develop meta-cognitive skills. 
\citet{qingsong2025raise} introduces a RL framework for dynamic instruction selection during fine-tuning.
In our work, we also value reflecting on reasoning processes, and we enhance metacognitive abilities through two-agent interaction and reinforcement learning at both end.

\subsection{RL in LLM}
\label{related_work:rl_in_llm}
Recent advancements in applying RL to LLMs have enhanced their reasoning and decision-making capabilities.
\citet{DBLP:journals/corr/abs-2503-20783} examines token-level optimization biases by introducing Dr. GRPO to stabilize policy gradients.
VAPO \citep{yue2025vapo} enhances PPO with value-aware perturbations and adaptive reward shaping to improve robustness in sparse-reward reasoning tasks.
DAPO \citep{yu2025dapo} provides a scalable, modular RL framework that integrates distributed rollout collection and dynamic replay buffers for reproducible training at scale.
SimpleRL-Zoo \citep{zeng2025simplerlzoo} conducts zero-shot RL experiments across open-base LLMs to uncover emergent cognitive behaviors under minimal reward signals.
Echo Chamber \citep{zhao2025echo} investigates how RL fine-tuning algorithms can amplify pretrained model biases and proposes regularization to mitigate over-amplification.
\citet{DBLP:journals/corr/abs-2405-15821} decomposes high-level language actions into token-level operations to achieve finer-grained credit assignment.
Some works push RL training for single-turn to multi-turn.
Search-R1 \citep{jin2025searchr1} trains LLMs to orchestrate multi-turn search strategies with RL-optimized decision policies to improve question-answering accuracy.
ArCHer \citep{zhou2024archer} employs a hierarchical, multi-turn RL architecture with manager and worker policies to efficiently handle long-horizon dialogue tasks.
RAGEN \citep{wang2025ragen} introduces trajectory filtering and critic modules within a multi-turn RL framework to stabilize learning and reduce shallow policy behaviors.



% Drawing from the bitter lesson \citep{sutton2019bitter}, two methods that appear to scale effectively are searching and learning, aligning with current trends in large language models \citep{xu2025towards}. At present, researchers are leveraging these methods to maximize the capabilities of individual transformers, while other efforts are exploring architectures that involve multiple interacting entities. In this paper, we examine this divergence within the context of LLM reasoning, a capability that allows large language models to solve problems through logical reasoning, step-by-step analysis, and inference \citep{wang2024openr}.

% \subsection{Single LLM Reasoning}

% Main research works in reasoning involving a single LLM utilize search-based and post-training methods. The fundamental elements of searching methods are text generation and evaluation. Generation schemes include In-Context Learning \citep{brown2020language}, Beam Search \citep{graves2012sequence}, and various tree-based searching \citep{snell2024scaling}; Evaluation approaches often use outcome accuracy, self-consistency \citep{wang2022self}, or process reward signal \citep{lightman2023let} as the criteria to select high-quality responses from the generated texts. Post-training method is another research line in opposition to pre-training. Popular training pipelines often involve specific data construction followed by Supervised Fine-tuning \citep{qin2024o1replicationjourneystrategic, ouyang2022training, hui2024qwen2, liu2024deepseek}, or reinforcement learning to interactively explore learning patterns \citep{wang2024openr, zhang2024llamaberrypairwiseoptimizationo1like, r1, xu2025towards}.

% \subsection{Multiple LLM Reasoning}
% Integrating multiple entities can potentially surpass the intelligence of the individual model \citep{chen2023agentverse}. With the rapid emergence of large language models showing a varying level of abilities, some studies have explored facilitating discussions among multiple off-the-shelf LLMs \citep{zhang2025if,chen2024s,wang2023unleashing,du2023improving,zhuge2023mindstorms,tang2023medagents,hao2025chatllm,akata2023playing,hong2023metagpt,zhang2024aflow}, taking the form of free discussion \citep{du2023improving,liang2023encouraging} or structured role assignments \citep{hong2023metagpt,zhang2024aflow}. Some have applied routing mechanisms to assign tasks to the most suitable expert models \citep{hu2024routerbench,stripelis2024tensoropera,ding2024hybrid,yue2025masrouter,chen2024routerdc} or merging mechanisms to develop more versatile models \citep{yadav2024matters,yu2024language,zhang2025nature}. Beyond aggregating static knowledge from multiple agents, multi-agent LLM training can also enhance reasoning capabilities. For example, multi-agent debates can generate diverse synthetic data, which can subsequently be used for supervised fine-tuning \citep{estornell2024accdebateactorcriticapproachmultiagent,li2024can,motwani2024maltimprovingreasoningmultiagent,dong2025stpselfplayllmtheorem,perez2022red,ye2025emergence,subramaniam2025multiagentfinetuningselfimprovement}. RL methods have also been adopted to improve LLM reasoning in areas such as alignment \citep{perez2022red,ma2023red} and legibility \citep{kirchner2024prover}. \cite{motwani2024maltimprovingreasoningmultiagent} utilize a three-agent system for generation and fine-tune the models using Direct Preference Optimization (DPO). Reinforcement Learning with Generative Reward Models (GenRM) \citep{mahan2024generative,ye2025iterative,jiao2024preference,wang2024self} represents another common approach of multi-agent training, where the reward signal is derived from the token probabilities of another LLM, coupled with the reasoning process. While our work aligns with these efforts, it diverges by using an additional tunable LLM to provide metacognitive instructions, guiding the low-level LLM during learning, rather than relying on a static GenRM. The most closely related works to ours are MAPoRL \citep{park2025maporlmultiagentpostcotrainingcollaborative} and COPRY \citep{DBLP:journals/corr/abs-2410-06101}. MAPoRL is a multi-agent debating framework that uses MARL with a learned verifier to fine-tune each LLM agent. COPRY duplicates an LLM into two agents, training them simultaneously in the roles of pioneer and observer using RL. \citet{shen2025satorireinforcementlearningchainofactionthought} trained with a novel Chain-of-Action-Thought (COAT) framework that embeds meta-action tokens for self-reflection and exploration into an autoregressive search process.
% However, unlike our approach, which explicitly separates metacognition from plan execution, these methods do not decompose the reasoning process but instead focus on improving direct chain-of-thought generation. Furthermore, our experiments are conducted on a larger scale and include more challenging problems.

% \subsection{Hierarchical Reasoning}\label{sec:ref3}

% Partitioning reasoning into hierarchical processes has been explored in prior research to make biological sense \citep{ye2018survey,langley2004hierarchical}. In the context of language models, a hierarchical structure has been used to facilitate diverse reasoning patterns, including planning \citep{puerta2025roadmap,sun2024retrieval,song2023llm,rana2023sayplan,chen2024autotamp,yan2023ask,xiao2024cellagent}, validation \citep{haji2024improvingllmreasoningmultiagent,xi2024enhancing} and self-refinement \citep{madaan2023self,kumar2024training,welleck2022generating}. For instance, EvalPlanner \citep{saha2025learningplanreason} is a framework that conducts reasoning through plan generation and execution. DOTS \citep{DBLP:journals/corr/abs-2410-03864} extends decomposition by integrating a tree-based searching method with Analysis, Solution, and Verification layers. Marco-o1 \citep{zhao2024marcoo1openreasoningmodels} focuses on open-ended problem-solving and abstract thinking, dynamically adjusting reasoning granularity and incorporating reflection mechanisms to enhance reasoning performance. Beyond these approaches, metacognition \citep{flavell1979metacognition} has been identified as another critical component of reasoning, referring to the intuitive understanding of one's own cognitive and reasoning processes \citep{gao2024meta,wang2023metacognitive}. \cite{wang2023metacognitive} proposed a metacognitive prompting strategy to improve large language model (LLM) capabilities. \cite{didolkar2024metacognitive} further developed a prompt-guided method that enables models to label math problems with the required skills and subsequently use these labels to solve new problems. \cite{gao2024meta} introduce meta-reasoner which use contextual multi-arm bandit to learn a high-level ``advisor'' over low-level reasoning process. \cite{xiang2025towards} provides a Meta-CoT framework to think about its own thinking. They use construction-based methods as well as reinforcement learning to develop meta-cognitive skills. 
% \citet{qingsong2025raise} introduces a RL framework for dynamic instruction selection during fine-tuning.
% In our work, we also value reflecting on reasoning processes, and we enhance metacognitive abilities through two-agent interaction and reinforcement learning at both end.

% \subsection{RL in LLM}
% Recent advancements in applying RL to LLMs have enhanced their reasoning and decision-making capabilities.
% \citet{DBLP:journals/corr/abs-2503-20783} examines token-level optimization biases by introducing Dr. GRPO to stabilize policy gradients.
% VAPO \citep{yue2025vapo} enhances PPO with value-aware perturbations and adaptive reward shaping to improve robustness in sparse-reward reasoning tasks.
% DAPO \citep{yu2025dapo} provides a scalable, modular RL framework that integrates distributed rollout collection and dynamic replay buffers for reproducible training at scale.
% SimpleRL-Zoo \citep{zeng2025simplerlzoo} conducts zero-shot RL experiments across open-base LLMs to uncover emergent cognitive behaviors under minimal reward signals.
% Echo Chamber \citep{zhao2025echo} investigates how RL fine-tuning algorithms can amplify pretrained model biases and proposes regularization to mitigate over-amplification.
% \citet{DBLP:journals/corr/abs-2405-15821} decomposes high-level language actions into token-level operations to achieve finer-grained credit assignment.
% Some works push RL training for single-turn to multi-turn.
% Search-R1 \citep{jin2025searchr1} trains LLMs to orchestrate multi-turn search strategies with RL-optimized decision policies to improve question-answering accuracy.
% ArCHer \citep{zhou2024archer} employs a hierarchical, multi-turn RL architecture with manager and worker policies to efficiently handle long-horizon dialogue tasks.
% RAGEN \citep{wang2025ragen} introduces trajectory filtering and critic modules within a multi-turn RL framework to stabilize learning and reduce shallow policy behaviors.

